\documentclass[10.5pt,a4paper]{article}
\usepackage[a4paper,margin=2.2cm,top=2.3cm,bottom=2.4cm]{geometry}
\usepackage{fontspec}
\usepackage{unicode-math}
\setmainfont{TeX Gyre Pagella}
\setsansfont{Inter}[Scale=0.88]
\setmonofont{DejaVu Sans Mono}[Scale=0.80]
\setmathfont{DejaVu Math TeX Gyre}[Scale=0.95]
\usepackage{microtype}
\usepackage{xcolor}
\definecolor{accent}{HTML}{1F4E79}
\definecolor{muted}{HTML}{5F6670}
\definecolor{codebg}{HTML}{F6F7F9}
\definecolor{rule}{HTML}{D9DCE1}
\usepackage{graphicx}
\usepackage{booktabs,tabularx,array}
\usepackage{enumitem}
\setlist{nosep,leftmargin=1.4em,topsep=3pt}
\usepackage[font=small,labelfont={bf,sf},labelsep=period,skip=4pt]{caption}
\usepackage{titlesec}
\titleformat{\section}{\sffamily\large\bfseries\color{accent}}{\thesection}{0.7em}{}
\titleformat{\subsection}{\sffamily\normalsize\bfseries}{\thesubsection}{0.6em}{}
\titleformat{\subsubsection}[runin]{\sffamily\small\bfseries}{}{0em}{}[.]
\titlespacing*{\section}{0pt}{14pt}{6pt}
\titlespacing*{\subsection}{0pt}{10pt}{4pt}
\usepackage{fancyhdr}
\pagestyle{fancy}\fancyhf{}
\renewcommand{\headrulewidth}{0.4pt}
\renewcommand{\headrule}{\hbox to\headwidth{\color{rule}\leaders\hrule height \headrulewidth\hfill}}
\fancyhead[L]{\sffamily\footnotesize\color{muted}\ReportShortTitle}
\fancyhead[R]{\sffamily\footnotesize\color{muted}S. K. Eshraghi Naeini}
\fancyfoot[C]{\sffamily\footnotesize\color{muted}\thepage}
\usepackage{listings}
\lstdefinelanguage{VerilogX}[]{Verilog}{morekeywords={localparam,integer,genvar,generate,endgenerate}}
\lstset{language=VerilogX,basicstyle=\ttfamily\footnotesize,keywordstyle=\color{accent}\bfseries,
  commentstyle=\color{muted}\itshape,stringstyle=\color{teal!70!black},numbers=left,numberstyle=\tiny\color{muted},
  numbersep=6pt,backgroundcolor=\color{codebg},frame=single,rulecolor=\color{rule},framesep=4pt,
  xleftmargin=14pt,breaklines=true,columns=fullflexible,keepspaces=true,showstringspaces=false,
  tabsize=4,aboveskip=6pt,belowskip=6pt}
\usepackage{tikz}
\usetikzlibrary{arrows.meta,positioning,shapes.geometric,calc,fit,automata}
\tikzset{blk/.style={draw=accent,thick,fill=accent!6,rounded corners=2pt,minimum height=8mm,minimum width=20mm,align=center,font=\sffamily\small},
         io/.style={font=\sffamily\footnotesize,align=center},
         arr/.style={-{Stealth[length=2mm]},semithick},
         st/.style={state,draw=accent,thick,fill=accent!6,minimum size=9mm,font=\sffamily\footnotesize,inner sep=1pt,align=center}}
\usepackage{xurl}
\usepackage[colorlinks=true,linkcolor=accent,urlcolor=accent,citecolor=accent]{hyperref}
\usepackage{siunitx}
\sisetup{group-separator={,},group-minimum-digits=4}
\newcommand{\sig}[1]{\texttt{#1}}
\newcommand{\pass}{\textcolor{green!45!black}{\textsf{\textbf{PASS}}}}
\newcommand{\repo}{https://github.com/seyedkiarash/seyedkiarash.github.io/tree/main/labs/logic-circuits-lab}
\newcommand{\note}[1]{\par\smallskip\noindent\colorbox{accent!7}{\parbox{\dimexpr\linewidth-2\fboxsep}{\small #1}}\par\smallskip}
\setlength{\parindent}{0pt}
\setlength{\parskip}{4pt}

\newcommand{\ReportTitle}[3]{%
  \thispagestyle{plain}
  {\sffamily\small\color{muted}Logic Circuits Laboratory \textbullet\ Department of Electrical Engineering, Sharif University of Technology\par}
  \vspace{2pt}
  {\sffamily\LARGE\bfseries #1\par}
  \vspace{1pt}
  {\sffamily\Large #2\par}
  \vspace{6pt}
  {\color{rule}\hrule height 0.6pt}
  \vspace{5pt}
  {\small
  \begin{tabular}{@{}l@{\hspace{2em}}l@{}}
    \textcolor{muted}{Student}  & Seyed Kiarash Eshraghi Naeini (403101383)\\
    \textcolor{muted}{Instructor} & Dr.\ Mohammad Reza Movahhedin \quad\textcolor{muted}{Teaching assistant:} Hossein Moghim\\
    \textcolor{muted}{Semester} & Fall 2025 \quad\textcolor{muted}{session:} #3\\
    \textcolor{muted}{Platform} & Zynq-7010 (XC7Z010-CLG400-1), \SI{50}{MHz} clock, Vivado 2019.1
  \end{tabular}\par}
  \vspace{4pt}
  {\color{rule}\hrule height 0.6pt}
  \vspace{8pt}}

\newcommand{\ReportShortTitle}{Experiment 3 -- Priority encoder}
\begin{document}
\ReportTitle{3}{Priority Encoder and Push-Button Polarity}{17 November and 8 December 2025}

\section*{Summary}
A 4-input priority encoder is implemented on the FPGA with the board push buttons as inputs and the
board LEDs as outputs. The first version treats the buttons as active-high; the second inverts them,
because the buttons of the board read 0 when pressed. Both were verified exhaustively for this report.
Code: \url{\repo/lab03-priority-encoder}.

\section{Design}
A priority encoder outputs the index of the highest-priority active input. Here the output is the button
number \emph{plus one}, so that 000 can mean ``no button pressed'':
\begin{table}[h]
\centering\small
\begin{tabular}{@{}ccccc@{}}
\toprule
\sig{button[3]} & \sig{button[2]} & \sig{button[1]} & \sig{button[0]} & \sig{LED[2:0]}\\
\midrule
1 & $\times$ & $\times$ & $\times$ & 100\\
0 & 1 & $\times$ & $\times$ & 011\\
0 & 0 & 1 & $\times$ & 010\\
0 & 0 & 0 & 1 & 001\\
0 & 0 & 0 & 0 & 000\\
\bottomrule
\end{tabular}
\end{table}
The priority is expressed directly by an \texttt{if}\ldots\texttt{else if} chain, which synthesis turns
into the equations $L_2=b_3$, $L_1=\overline{b_3}(b_2+b_1)$, $L_0=\overline{b_3}(b_2+\overline{b_1}b_0)$
(2 LUTs).

\begin{lstlisting}
always @(*) begin
    if      (A[3]) c = 3'b100;
    else if (A[2]) c = 3'b011;
    else if (A[1]) c = 3'b010;
    else if (A[0]) c = 3'b001;
    else           c = 3'b000;
end
\end{lstlisting}

\section{Part 3 vs.\ Part 5 -- button polarity}
In Part 3 the buttons were connected directly. On the board, however, a button reads \emph{0} when it is
pressed, so with no button pressed all inputs are 1 and the encoder shows 100. Part 5 places the same
encoder (\sig{encoderpr}) behind an inverter, \sig{.A(\textasciitilde button)}, after which the LEDs
behave as intended. The same inversion is used in the later experiments wherever the board buttons
serve as data inputs.

\section{Verification}
\begin{table}[h]
\centering\small
\begin{tabularx}{\linewidth}{@{}lXr@{}}
\toprule
Testbench & What is checked & Result\\
\midrule
\sig{tb\_priority\_encoder} & all 16 input combinations, active-high & \pass\\
\sig{tb\_part5} & all 16 combinations, active-low buttons & \pass\\
\bottomrule
\end{tabularx}
\end{table}

\section{Conclusion}
The priority encoder is a direct translation of its truth table into a prioritised conditional chain.
The second version is a reminder to check the electrical polarity of every input before testing logic on
the board.

\end{document}
